\begin{ZhChapter}

\chapter{緒論}

\section{研究背景與動機}

慢性阻塞性肺病（Chronic Obstructive Pulmonary Disease, COPD）為一種與長期氣道阻塞、肺實質破壞及肺功能下降相關之慢性呼吸道疾病。臨床診斷主要依賴肺功能檢查，然而胸部電腦斷層影像（computed tomography, CT）能呈現肺實質、氣道與血管等結構變化，因此具有作為輔助診斷與疾病表徵分析之潛力。

\section{研究目的}

本研究旨在建立以胸部 CT 影像為基礎之 COPD 輔助診斷流程，並從可解釋量化特徵與三維影像模型兩個方向進行分析。前者擷取肺氣腫、氣道、血管與肺容積等量化生物標記，後者使用三維 CT 模型學習肺功能相關影像表徵。

\section{研究範圍與貢獻}

本研究範圍聚焦於胸部 CT 影像於 COPD 輔助診斷之應用，並不取代臨床肺功能檢查或醫師診斷。本研究所討論之分類與分級結果，主要作為影像輔助分析與研究性評估依據。主要貢獻如下：

\begin{enumerate}
    \item 建立胸部 CT 影像之 COPD 量化特徵分析流程，並以 12 項影像生物標記進行 Normal/Abnormal 分類。
    \item 評估三維 CT 深度學習模型於塌陷角角度分類與 GOLD 分級任務之可行性。
    \item 探討小樣本與類別不平衡情境下，資料增強與平衡取樣對模型效能之影響。
    \item 比較量化特徵方法與三維影像模型於 COPD 輔助診斷中之互補性。
\end{enumerate}

\section{論文架構}

本論文共分為五章。第二章整理相關背景與文獻，第三章說明研究材料與方法，第四章呈現實驗結果與討論，第五章總結研究成果、限制與未來方向。

\end{ZhChapter}
